\clearpage
\begin{summarybox}
\textbf{\textcolor{CSummary}{一句话核心：两条出自特定点的弦，斜率互为相反数；第三条弦恒与 $y=-x$ 平行。}}

\[
 \boxed{y^2=2px\ (p>0),\qquad A=\left(\frac p2,p\right)}
\]
\[
 \boxed{k_{AM}=-k_{AN}\quad\Longrightarrow\quad k_{MN}=-1}.
\]
\textbf{最快的代数抓手：}令 $M=P(u)$、$N=P(v)$、$P(t)=(\frac{t^2}{2p},t)$，则
\[
 \boxed{k_{P(u)P(v)}=\frac{2p}{u+v}},\qquad
 \boxed{u+v=-2p}.
\]
\[
 \boxed{\frac1{k_{MN}}=\frac1{k_{AM}}+\frac1{k_{AN}}-1=-1}.
\]
\textbf{三秒识别流程：}
\[
 \boxed{A=(p/2,p)}\ \Longrightarrow\
 \boxed{k_{AM}=-k_{AN}}\ \Longrightarrow\
 \boxed{u+v=-2p}\ \Longrightarrow\
 \boxed{k_{MN}=-1}.
\]
\begin{itemize}[leftmargin=*,itemsep=0.36em]
 \item \textbf{最低门槛：}直接算 $k_{AM},k_{AN}$，消元得 $u+v=-2p$。
 \item \textbf{最快：}记住\textbf{斜率倒数恒等式}，代入相反数立刻得 $-1$。
 \item \textbf{常规：}令两直线斜率为 $t,-t$，与抛物线联立，得到两点纵坐标之和。
 \item \textbf{边界：}三点互异；$AM,AN$ 斜率须存在；不要把“斜率相反”误当成“纵坐标相反”。
 \item \textbf{推广：}若定点纵坐标为 $a\ne0$，结论改为 $k_{MN}=-p/a$。
\end{itemize}
\textbf{\textcolor{CSummary}{可迁移的思路：}}\textbf{求两动点连线的斜率，先看纵坐标之和是否被已知条件固定。}本题不需要分别求出两动点。
\end{summarybox}


\clearpage
\begin{summarybox}
\textbf{补充速记：C097 $\longrightarrow$ C098}
\[
 \boxed{\frac1{k_{MN}}=\frac1{k_{AM}}+\frac1{k_{AN}}-\frac2{k_{OA}}}
\]
\[
 k_{AM}=-k_{AN}\quad\Longrightarrow\quad
 \boxed{k_{MN}=-\frac12k_{OA}}.
\]
\textbf{一般推广：}$A$ 可换成抛物线上任意非顶点定点 $A_0=(\frac{a^2}{2p},a)$，此时
\[
 \boxed{k_{MN}=-\frac12k_{OA_0}=-\frac pa}\qquad(a\ne0).
\]
原题 $k_{OA}=2$，所以 $k_{MN}=-1$。使用 C097 的倒数关系时，若 $M$ 或 $N$ 是顶点 $O$，改用两点斜率公式验证。
\end{summarybox}
